\documentclass[11pt]{article}
\usepackage[a4paper,margin=0.72in,top=0.72in,bottom=0.78in]{geometry}
\usepackage[T1]{fontenc}
\usepackage{charter}
\usepackage[scaled=0.92]{helvet}
\usepackage{amsmath,amssymb}
\usepackage{booktabs}
\usepackage{titlesec}
\usepackage{xcolor}
\usepackage{enumitem}
\usepackage[hidelinks]{hyperref}

\newcommand{\num}[1]{#1}
% Generated by src/study.py - do not edit by hand.
% Rerun 'python3 run.py' to refresh these numbers.
\newcommand{\NLocations}{32}
\newcommand{\NComplaintSites}{31}
\newcommand{\NDorms}{21}
\newcommand{\NNewCampus}{24}
\newcommand{\NOldCampus}{8}
\newcommand{\NTunnelMetres}{94.3}
\newcommand{\NMeanWalk}{5.7}
\newcommand{\NMeanSameCampus}{2.1}
\newcommand{\NMeanCrossCampus}{11.3}
\newcommand{\NMaxWalk}{16.1}
\newcommand{\NFifoObj}{3602.6}
\newcommand{\NFifoTard}{3241.0}
\newcommand{\NFifoTravel}{26.8}
\newcommand{\NFifoSla}{85.1}
\newcommand{\NFifoUrgent}{37.0}
\newcommand{\NFifoUrgentWait}{4.47}
\newcommand{\NFifoUrgentNinety}{7.26}
\newcommand{\NFifoBacklog}{6.1}
\newcommand{\NFifoClosed}{1657}
\newcommand{\NFifoUrgentLate}{153}
\newcommand{\NFifoUrgentUnavoidable}{32}
\newcommand{\NFifoUrgentAvoidable}{121}
\newcommand{\NSavObj}{2766.2}
\newcommand{\NSavTard}{2563.8}
\newcommand{\NSavTravel}{15.0}
\newcommand{\NSavSla}{90.5}
\newcommand{\NSavUrgent}{55.5}
\newcommand{\NSavUrgentWait}{3.47}
\newcommand{\NSavUrgentNinety}{6.33}
\newcommand{\NSavBacklog}{3.3}
\newcommand{\NSavClosed}{1665}
\newcommand{\NSavUrgentLate}{109}
\newcommand{\NSavUrgentUnavoidable}{32}
\newcommand{\NSavUrgentAvoidable}{77}
\newcommand{\NPriObj}{358.0}
\newcommand{\NPriTard}{106.5}
\newcommand{\NPriTravel}{18.6}
\newcommand{\NPriSla}{97.9}
\newcommand{\NPriUrgent}{85.7}
\newcommand{\NPriUrgentWait}{2.92}
\newcommand{\NPriUrgentNinety}{4.07}
\newcommand{\NPriBacklog}{0.2}
\newcommand{\NPriClosed}{1665}
\newcommand{\NPriUrgentLate}{35}
\newcommand{\NPriUrgentUnavoidable}{32}
\newcommand{\NPriUrgentAvoidable}{3}
\newcommand{\NImpObj}{319.9}
\newcommand{\NImpTard}{108.3}
\newcommand{\NImpTravel}{15.7}
\newcommand{\NImpSla}{98.0}
\newcommand{\NImpUrgent}{86.9}
\newcommand{\NImpUrgentWait}{2.42}
\newcommand{\NImpUrgentNinety}{4.05}
\newcommand{\NImpBacklog}{0.6}
\newcommand{\NImpClosed}{1665}
\newcommand{\NImpUrgentLate}{32}
\newcommand{\NImpUrgentUnavoidable}{32}
\newcommand{\NImpUrgentAvoidable}{0}
\newcommand{\NUrgentTotal}{245}
\newcommand{\NUrgentImpossible}{32}
\newcommand{\NTravelCutPct}{41.4}
\newcommand{\NObjCutPct}{91.1}
\newcommand{\NWalkSavedPerDay}{100.0}
\newcommand{\NBbMaxTickets}{10}
\newcommand{\NBbGapWorst}{0.0}
\newcommand{\NBbServed}{7}
\newcommand{\NBbNodesArbCold}{22,721}
\newcommand{\NBbNodesArbWarm}{10,131}
\newcommand{\NBbOrderCutPct}{52.9}
\newcommand{\NBbWarmCutPct}{5.3}
\newcommand{\NBbNodesWarm}{10,131}
\newcommand{\NBbNodesCold}{10,700}
\newcommand{\NBbNodeCutPct}{5.0}
\newcommand{\NBbSeconds}{0.02}
\newcommand{\NBbSavWorstGap}{4342.5}
\newcommand{\NBbSavWorstJobs}{8}
\newcommand{\NBbSavWorstTimes}{44.4}
\newcommand{\NCrewSixObj}{2717.6}
\newcommand{\NCrewSixSla}{90.3}
\newcommand{\NCrewSixUrgent}{79.6}
\newcommand{\NCrewNineObj}{319.9}
\newcommand{\NCrewNineSla}{98.0}
\newcommand{\NCrewNineUrgent}{86.9}
\newcommand{\NCrewTwelveObj}{293.1}
\newcommand{\NCrewTwelveSla}{98.1}
\newcommand{\NCrewTwelveUrgent}{86.9}
\newcommand{\NCrewRebalObj}{2355.7}
\newcommand{\NCrewRebalSla}{93.4}
\newcommand{\NCrewRebalUrgent}{85.3}
\newcommand{\NSurgeLowFifoSla}{85.1}
\newcommand{\NSurgeLowFifoBacklog}{6.1}
\newcommand{\NSurgeLowFifoUrgent}{37.0}
\newcommand{\NSurgeLowImpSla}{98.0}
\newcommand{\NSurgeLowImpBacklog}{0.6}
\newcommand{\NSurgeLowImpUrgent}{86.9}
\newcommand{\NSurgeHighFifoSla}{53.9}
\newcommand{\NSurgeHighFifoBacklog}{113.5}
\newcommand{\NSurgeHighFifoUrgent}{18.6}
\newcommand{\NSurgeHighImpSla}{88.1}
\newcommand{\NSurgeHighImpBacklog}{23.8}
\newcommand{\NSurgeHighImpUrgent}{83.9}
\newcommand{\NSeedCount}{5}
\newcommand{\NSeedImpMin}{277.9}
\newcommand{\NSeedImpMax}{348.5}
\newcommand{\NSeedFifoMin}{3562.4}
\newcommand{\NSeedFifoMax}{4825.5}
\newcommand{\NDayTickets}{58}
\newcommand{\NDayFifoTravel}{239.0}
\newcommand{\NDayImpTravel}{160.9}
\newcommand{\NDays}{30}
\newcommand{\NTicketsPerDay}{58}
\newcommand{\NWorkingMinutes}{450}
\newcommand{\NAlpha}{0.05}
\newcommand{\NAlphaZeroTravel}{18.3}
\newcommand{\NAlphaWorthMinutes}{2.6}


\definecolor{ink}{HTML}{1A1A1A}
\definecolor{accent}{HTML}{1F4E79}
\color{ink}

\titleformat{\section}{\normalfont\bfseries\sffamily\color{accent}\large}{\thesection.}{0.6em}{}
\titlespacing{\section}{0pt}{10pt}{3pt}
\titleformat{\subsection}{\normalfont\bfseries\sffamily\color{ink}\normalsize}{}{0em}{}
\titlespacing{\subsection}{0pt}{7pt}{1.5pt}

\setlist[itemize]{leftmargin=1.2em,itemsep=1.5pt,topsep=2pt}
\setlength{\parskip}{2.5pt}
\setlength{\parindent}{0pt}
\linespread{1.0}

\begin{document}

\begin{center}
{\sffamily\bfseries\Large How the dorm maintenance dispatch model works}\\[4pt]
{\sffamily\normalsize A plain explanation of the data, the methods and the findings}\\[3pt]
{\small Operations Management \& Logistics group project}
\end{center}
\vspace{4pt}
\hrule
\vspace{8pt}

This is the companion to the main report. The report argues the case; this
one explains the machinery. Every number quoted here comes straight from the
code's own output, so it cannot drift from what the model does.

\section{The problem, in plain words}

The Estate Office gets maintenance complaints all day: a dead tubelight here,
a dripping tap there, a window that will not shut. Each goes to a technician
who walks to the room, does the repair, and moves on. Right now the queue is
worked roughly in the order complaints arrive, which sounds fair but wastes a
great deal of walking and ignores urgency completely. The question the model
answers is narrow: \emph{given the complaints waiting at the start of a day,
which technician should take which jobs, and in what order?}

Two features of IIMA make this harder than it sounds. First, a main road cuts
the campus in two and there is no crossing at road level, so every trip
between the halves goes through the walkway tunnel and costs about
\NMeanCrossCampus{} minutes against \NMeanSameCampus{} minutes for a trip
that stays on one side. Second, the crews are not in one place. The
electricians work out of a store room in Dorm 7 on the main campus, and the
plumbers and carpenters share one in Dorm 24 on the new campus. The two are
only about \num{460}~metres apart in a straight line but \num{10.6} minutes
apart on foot. So each trade has its own far side of the campus, and the same
job is cheap for one trade and expensive for another.

\section{What is real and what is invented}

This distinction matters, because it is the first thing anyone will ask.

\textbf{Real:} the campus. Coordinates for all \NLocations{} locations come
from OpenStreetMap, including the walkway tunnel, all \NTunnelMetres{} metres
of it; \NNewCampus{} are on the new campus and \NOldCampus{} on the main
campus. The exceptions are Dorm 7 and Dorm 24, which OpenStreetMap does not
label individually. Those were placed from the official IIMA campus map,
which groups the dormitories as D1--D18 on the main campus and D19--D27,
D28--D35 and D36--D43 on the new one, each on the footprint of a real
building in the right group.

\textbf{Invented:} the complaints. Nobody was going to hand us the Estate
Office's ticket log, so the model generates its own. The point was never to
guess the true numbers but to put every assumption in one file where it can
be argued with. Give it real figures in \texttt{src/instance.py} and rerun,
and the report rewrites its own numbers.

\section{Every assumption, in one place}

\begin{center}\small
\begin{tabular}{@{}llp{7.4cm}@{}}
\toprule
\textbf{What} & \textbf{Value} & \textbf{Why that value} \\
\midrule
Working day & \NWorkingMinutes{} min & 09:00--17:30 with an hour for lunch \\
Complaints per day & \NTicketsPerDay{} & An average; the actual count each day is drawn at random \\
Share from dorms & 85\% & Dorms have far more people in them than offices \\
Trade mix & 45 / 35 / 20\% & Electrical / plumbing / carpentry. Lights, fans and sockets dominate \\
Repair time & 22 / 38 / 30 min & Median per trade; plumbing takes longest \\
Priority split & 15 / 55 / 30\% & Urgent / normal / minor \\
Lateness weight & 10 / 4 / 1 & How badly we mind an hour late in each class \\
Promised response & 4 h / 1 day / 3 days & The deadline per priority class \\
Access windows & 40 / 30 / 30\% & Anytime / mornings only / afternoons only, for dorm rooms \\
Walking speed & 75 m/min & An unhurried pace carrying a toolbag \\
Detour factor & \num{1.30} & You walk around lawns and buildings, not through them \\
Tunnel penalty & 2 min & Down the ramp, along, and back up \\
Price of walking & $\alpha=\NAlpha$ & Per minute, against one weighted hour of lateness \\
Crew & 3 per trade & Nine technicians in total \\
\bottomrule
\end{tabular}
\end{center}

Two of these carry more weight than the rest. The \textbf{access window mix}
turns out to decide the final result, and it is the one assumption most worth
replacing with real data. The \textbf{price of walking} is deliberately tiny:
walking is there to break ties, not to drive the answer. We never accept a
later repair in order to save a walk.

\section{What the model is trying to do}

It minimises
\[
\underbrace{\sum_i w_i \times (\text{hours job } i \text{ was late})}_{\text{the thing we care about}}
\;+\;
\underbrace{\alpha \times (\text{total walking minutes})}_{\text{the tiebreaker}}
\]
subject to four rules: every job is done at most once and only by a
technician of the matching trade; nobody starts a job before its access
window opens or after it shuts; nobody is still working after the shift ends;
and everyone is back at their own trade's store room by the end of the day.

One subtlety is a natural question. A job that does not fit today is not
lost: it stays in the queue and reappears tomorrow. But if deferring were
free, the cheapest plan would be to do nothing. So whenever the model
considers leaving a job, it charges that job the lateness it would collect if
finished at the end of tomorrow. A minor complaint with three days left costs
nothing to postpone, which is correct; an urgent one costs a great deal.

\section{The methods, from crudest to best}

Four methods were built deliberately in order, so each one exposes what the
previous one was missing.

\subsection{1. First come, first served -- the baseline}
Take complaints in the order raised, give each to the technician of the right
trade with the least work so far. No attention to geography or deadlines.
This is roughly what happens today, and it is what everything else is
measured against.

\subsection{2. Clarke and Wright's savings heuristic -- fixing the walking}
Start from the worst sensible plan: every job is its own trip out from the
store room and straight back. Then ask what you save by doing jobs $i$ and
$j$ on one trip instead of two:
\[
s_{ij} \;=\; c_{0i} + c_{0j} - c_{ij},
\]
where $0$ is the store room of the trade doing the work. Compute this for
every pair, sort from biggest saving down, and work through the list merging
routes whenever the merge is legal -- the two jobs must be on different
routes, each still at the end of its route, and the merged route must still
fit the shift and respect every window.

Why it suits this campus is the nicest thing in the project. For an
electrician every new-campus job has a large $c_{0i}$, being on the far side
of the tunnel from Dorm 7, so $s_{ij}$ is large for any pair of new-campus
jobs and the savings list is dominated by merges that bundle them into one
crossing. For a plumber working out of Dorm 24 the main-campus jobs behave
that way instead. Each trade batches its own far-side work, in opposite
directions, and neither was ever told the tunnel exists: the heuristic finds
it from the travel times alone.

\subsection{3. Priority insertion -- fixing the urgency}
Savings is blind to deadlines, so this method takes the opposite view. Insert
jobs one at a time; at each step, for every unplaced job, work out the
cheapest place it could slot into anyone's route, costing $\Delta$ extra
minutes. Then pick the job with the best ratio of what you avoid by doing it
today to what it costs to fit in. This is the ``bang per buck'' rule, and it
behaves sensibly when there is more work than time: the urgent jobs go in
first and the minor ones are what gets left over.

\subsection{4. Savings plus improvement -- the one we recommend}
Run savings to get sensible routes, then improve them with three moves,
repeated until nothing gets better:
\begin{itemize}
\item \textbf{2-opt:} reverse a stretch of a route. On a map this removes
  crossings -- if a route doubles back on itself, flipping the segment
  between the crossing legs untangles it.
\item \textbf{Or-opt within a route:} lift out a run of one to three
  consecutive jobs and reinsert it somewhere else in the same route. Gentler
  than 2-opt, because it leaves the order of everything else alone, which
  matters when windows are tight.
\item \textbf{Or-opt across routes:} move a short run of jobs from one
  technician to another of the same trade. This is the move that matters most
  in practice. Savings balances routes by geography but not by workload, so
  one technician ends up with a packed day while a colleague has a free
  afternoon; shifting a couple of jobs across fixes it.
\end{itemize}
Because access windows can make any of these illegal, every candidate move is
re-checked for feasibility rather than assumed valid.

\subsection{Checking the answer: branch and bound}
A heuristic comes with no guarantee, so for small instances we find the true
optimum and compare. The search builds a tree where each node is a partial
route and each child adds one more job. At every node it computes a lower
bound on the best any completion could achieve; if that bound is already
worse than the best plan found so far, the whole branch is thrown away
without being explored.

The bound has two parts. Travel: what has been walked plus the walk back to
the store room, valid because the times come from a shortest-path
calculation, so going anywhere else first can only cost more. Lateness: each
job still open is charged the cheaper of its two fates, the least lateness it
could pick up if done today or the deferral charge if left. Both ignore
ordering, which is why they never overstate the true cost.

\textbf{The result:} savings plus improvement found the exact optimum on
every instance up to \NBbMaxTickets{} jobs, a worst gap of \NBbGapWorst\%.
It also defers exactly the jobs the optimum defers. Two details are worth
quoting if asked: exploring the nearest job first cuts the search from
\NBbNodesArbCold{} nodes to \NBbNodesCold{}, a \NBbOrderCutPct\% saving,
while starting from the heuristic's answer saves only a further
\NBbWarmCutPct\%. The warm start matters less than you would expect, because
once stopping early is a legal plan the search finds a decent incumbent
within a few nodes anyway.

\section{What came out of it}

Over \NDays{} simulated days, against the first-come-first-served baseline:

\begin{center}\small
\begin{tabular}{@{}lrrl@{}}
\toprule
& \textbf{Baseline} & \textbf{Optimised} & \\
\midrule
Walking per technician per day & \NFifoTravel{} min & \NImpTravel{} min & down \NTravelCutPct\% \\
All complaints closed on time & \NFifoSla\% & \NImpSla\% & \\
Urgent complaints closed on time & \NFifoUrgent\% & \NImpUrgent\% & \\
Average queue left at day's end & \NFifoBacklog{} & \NImpBacklog{} & \\
Avoidable late urgent complaints & \NFifoUrgentAvoidable{} & \NImpUrgentAvoidable{} & in 30 days \\
\bottomrule
\end{tabular}
\end{center}

\subsection{The finding worth leading with}
Of the \NUrgentTotal{} urgent complaints raised over the month,
\NUrgentImpossible{} could never have met the four-hour promise no matter who
was sent or in what order: the room can only be opened in the afternoon and
the repair itself runs past the deadline. Even a technician already standing
at the door could not have made it.

Now count the misses. The optimised plan closes \NImpUrgentLate{} urgent
complaints late, and all \NImpUrgentUnavoidable{} of them are of exactly that
kind. \emph{The two numbers are the same number.} Every urgent complaint that
could have been closed on time was. The scheduling has run out of room to
improve, and what is left is a problem with the promise, not with the plan.
Fixing it is a policy decision -- either stop making a four-hour promise on
complaints that can only be accessed in the afternoon, or give the crew a way
into the room without the occupant, such as a dorm staff member with a key.
The second is worth about \NUrgentImpossible{} rescued complaints a month.

\subsection{Other things we learned}
\begin{itemize}
\item \textbf{Do not hire.} Six technicians collapses (objective
  \NCrewSixObj{} against \NCrewNineObj{} for nine). Nine is enough. Twelve
  improves it only to \NCrewTwelveObj{}, which buys almost nothing.
\item \textbf{Do not rebalance the trades to match average workload.} This
  one went against our own guess. Moving to four electricians, four plumbers
  and one carpenter -- the same nine people, allocated to match average
  demand -- gives \NCrewRebalObj{}, far worse than \NCrewNineObj{}. The quiet
  trade is quiet \emph{on average} and not on a bad day, and a single
  carpenter who falls behind never catches up.
\item \textbf{It holds up under pressure.} At a surge of 85 complaints a day,
  the baseline collapses to \NSurgeHighFifoSla\% on time with a backlog of
  \NSurgeHighFifoBacklog{}, while the optimised plan holds \NSurgeHighImpSla\%
  with a backlog of \NSurgeHighImpBacklog{}.
\item \textbf{It is not a fluke of one random seed.} Across \NSeedCount{}
  independent ticket streams the optimised objective ranges
  \NSeedImpMin--\NSeedImpMax{} and the baseline \NSeedFifoMin--\NSeedFifoMax{}.
  The ranking never changes.
\item \textbf{Most of the walking saving is a by-product.} Set $\alpha$ to
  zero, telling the model not to care about walking at all, and walking rises
  only from \NImpTravel{} to \NAlphaZeroTravel{} minutes. Planning the day
  sensibly saves the walking almost by itself; pricing it explicitly buys the
  last \NAlphaWorthMinutes{} minutes.
\end{itemize}

\section{Questions you should be ready for}

\textbf{``Your complaint data is made up, so why believe any of this?''}
Because the comparison is what matters, not the level. Both methods face the
same complaints, the same campus and the same crew, so the gap between them
is real even if the absolute numbers are not. And the ranking is stable
across five independent streams and across every assumption we varied.

\textbf{``Isn't 2-opt going to break the access windows?''} Yes, often. That
is why every candidate move is re-evaluated for feasibility instead of being
assumed valid. It is slower, and it is the only correct way to do it once
jobs have windows.

\textbf{``How do you know the heuristic is any good?''} On every instance
small enough to solve exactly, it matched the optimum. Not proof for a
sixty-job day, but the strongest evidence available without a commercial
solver.

\textbf{``What are its limits?''} It plans in the morning and never revises,
so an urgent complaint raised at eleven waits for tomorrow. Repair times are
assumed known in advance. And the exact check covers one technician with ten
jobs, not nine with sixty.

\end{document}
